\begin{resumo}
O presente documento formaliza a Arquitetura de Software do SIGEA-GTT (Sistema Inteligente de Gestão de Eventos Adversos -- \textit{Global Trigger Tool}), plataforma computacional concebida para a auditoria retrospectiva de prontuários clínicos, formação médica/enfermagem e vigilância epidemiológica hospitalar, desenvolvida em parceria entre o Departamento de Computação (DCOMP) e o Departamento de Enfermagem da Universidade Federal de Sergipe (UFS). Estruturado com base no consagrado Modelo de Visões Arquiteturais 4+1 de Philippe Kruchten, este trabalho descreve a organização do sistema a partir de cinco perspectivas complementares: (i) a \textbf{Visão de Cenários (+1)}, que ancora as decisões arquiteturais nos casos de uso críticos de autenticação institucional restrita, auditoria sob cronômetro regressivo estrito de 20 minutos, aplicação integrada de seis Ferramentas da Qualidade e cálculo automatizado de indicadores epidemiológicos do IHI ($TI_{1000}$, $TI_{100}$, $P_{EA}$); (ii) a \textbf{Visão Lógica}, detalhando a decomposição em camadas orientada a \textit{Clean Architecture} e \textit{Domain-Driven Design} (DDD) no backend Spring Boot 3.3 (Java 21) e a arquitetura modular reativa baseada em \textit{Standalone Components} e \textit{Signals} no frontend Angular 18, acompanhadas dos diagramas de classes de domínio e entidade-relacionamento (DER); (iii) a \textbf{Visão de Processos}, abordando o modelo de concorrência, o gerenciamento de conexões via HikariCP, o controle transacional ACID, a temporização cliente-side reativa e diagramas de sequência UML em TikZ para os fluxos operacionais centrais; (iv) a \textbf{Visão de Desenvolvimento}, caracterizando a organização do monorepo, a estrutura de pacotes Maven e Angular, o diagrama de componentes de software e o \textit{Quality Gate} inegociável de 100\% de cobertura de testes automatizados e 100\% de documentação de código; e (v) a \textbf{Visão Física / Implantação}, que delimita a topologia de nós, contêineres e comunicação HTTPS/TLS, garantindo isolamento total em relação aos sistemas do Hospital Universitário (HU-UFS/EBSERH) em estrita conformidade com a Lei Geral de Proteção de Dados (LGPD). O documento consolida ainda a matriz de rastreabilidade entre as visões arquiteturais e os 10 Requisitos Não Funcionais da norma ISO/IEC 25010.

\vspace{\onelineskip}
\noindent
\textbf{Palavras-chave}: Arquitetura de Software. Modelo 4+1 de Kruchten. Global Trigger Tool. Clean Architecture. Segurança do Paciente. Universidade Federal de Sergipe.
\end{resumo}
